\ \\[-10cm]
La parte anterior de esta guía exploró algunos de los principales retos de comunicar varios ordenadores entre sí.
\concept{Internet} es una solución que funciona para esos retos: la que los humanos hemos creado a lo largo del último medio siglo,
y que hoy resulta central en muchas actividades sociales, comerciales y políticas. 
% 
Esta parte estudia las principales prestaciones que ofrece Internet, así como el \conceptRef{suit}{conjunto} de \conceptRef{protocol}{protocolos} 
que se emplean para conseguirlas.


\chapterimage{earth.jpg}

\chapter{Internet}\label{sec:internet}

Las comunicaciones modernas se basan en \conceptRef{stack}{pilas} de red organizadas por capas. 
Quizá acabas de descargar esta guía de Internet.
% 
¿Qué capas se usan exactamente? ¿Qué protocolos se emplean, y con qué finalidad?
% 
Este capítulo ofrece una primera aproximación para responder a estas preguntas.\\[-1cm]

\section{La pila TCP/IP}


\concept{Internet} es una solución práctica al problema de la comunicación global entre ordenadores,
cuyo desarrollo comenzó a finales de los años 60 -- principios de los 70, y todavía continúa.

La conmutación de \conceptRef{packet}{paquetes} se concibió a principios de los años 60, y la primera 
implementación de área extensa, \concept{ARPANET}, comenzó en 1969. 
% 
Un reto clave en aquel momento era interconectar las \conceptRef{network}{redes} existentes, 
propiedad de organizaciones distintas y basadas en tecnologías y \conceptRef{protocol}{protocolos} diferentes.
% 
Para que aquello fuese viable, había que cumplir todos los requisitos siguientes:
\begin{itemize}
    \item Las redes existentes debían funcionar sin actualizar la infraestructura física 
      ni sus protocolos, solo la forma en que se usan esos protocolos.
    \item Las redes existentes debían seguir gestionándose de forma independiente a nivel local (organización interna autónoma).
    \item Los paquetes no siempre se entregan: algunos se pierden.
    \item Los mensajes que llegan pueden no estar en orden, y pueden repetirse.
\end{itemize}

La solución adoptada por ARPANET --y más tarde por el mundo entero-- 
es el conjunto de protocolos \concept{TCP/IP}, descrito por primera vez en 1974. 
Es un enfoque por capas, diseñado para ``enchufarse'' encima de 
las \conceptRef{LAN}{LANs} (redes de área local) existentes, con el objetivo de ofrecer 
capacidades de red a todas las aplicaciones de usuario.

\begin{center}
\begin{bytefield}{16}
\begin{rightwordgroup}[curlyshrinkage=10pt]{User programs}
\wordbox{3}{Layer 7: Application}
\end{rightwordgroup} \\
\begin{rightwordgroup}[curlystyle=\color{color1},curlyshrinkage=10pt]{{\color{color1} TCP/IP}}
\wordbox{1}{Layer 4: Transport}\\
\wordbox{1}{Layer 3: Internet}
\end{rightwordgroup} \\
\begin{rightwordgroup}[curlyshrinkage=10pt]{Existing/future LAN technologies}
\wordbox{1}{Layer 2: Network}\\
\wordbox{1}{Layer 1: Physical}
\end{rightwordgroup}
\end{bytefield}
\end{center} 
 
En resumen, los principales servicios que ofrecen estas capas son los siguientes:
\begin{enumerate}
\item[\textbf{Capa 1}] \textbf{(Física)}: 
  Enviar datos \conceptRef{digital}{digitales} a través de medios \conceptRef{analog}{analógicos}.
\item[\textbf{Capa 2}] \textbf{(Red)}:
  Identificar y comunicar dispositivos dentro de una misma \concept{LAN}.\\
  A veces se usa el término capa de ``Enlace de datos'' (del modelo OSI).
\item[\textbf{Capa 3}] \textbf{(Internet)}:
  Identificar y comunicar dispositivos entre todas las \conceptRef{LAN}{LANs}.\\
  A veces se usa el término capa de ``Red'' (también de OSI).
\item[\textbf{Capa 4}] \textbf{(Transporte)}:
  Identificar y comunicar programas (servicios).\\
  Además (elige \textit{solo} una):
  \begin{enumerate}[label=\alph*)]
   \item Asegurar que un \conceptRef{stream}{flujo} de datos llega correctamente al otro extremo, 
   y en orden \\-- \textit{pero} se introduce un \concept{overhead} (sobrecarga) 
   y hay que establecer \conceptRef{connection}{conexiones} (\concept{TCP}).
     
   \item Usar mensajes sencillos y relativamente eficientes, con poco overhead\\
   -- \textit{pero} los mensajes pueden perderse, reordenarse o duplicarse (\concept{UDP}).
  \end{enumerate}
\item[\textbf{Capa 7}] \textbf{(Aplicación)}:
  Hacer cualquier cosa que el usuario pueda necesitar. \\
  No olvides la seguridad (\eg, criptografía) 
  ni la eficiencia (\eg, compresión): la capa 4 y las inferiores no proporcionan eso.
\end{enumerate}

\section{Visión general de los protocolos TCP/IP}

Se definen múltiples \conceptRef{protocol}{protocolos} para implementar la funcionalidad
de cada \conceptRef{layer}{capa}. Estas definiciones incluyen el \conceptRef{packet format}{formato de paquete} necesario,
las reglas para intercambiar esos paquetes y su semántica.
% 
Estos protocolos deben entenderse como un sistema en el que cada parte tiene tareas muy concretas.

En la figura siguiente se muestran algunos de los protocolos más relevantes para 
las comunicaciones a través de Internet. En ella:
\begin{itemize}
\item Las flechas continuas indican que un protocolo se encapsula dentro de otro. 
  Si el destino es una capa, debe usarse alguno de los protocolos de esa capa.

\item Las cajas discontinuas indican protocolos auxiliares definidos para ayudar a que otros funcionen. 
Las flechas discontinuas indican qué capa usa qué protocolo auxiliar.

\item Las cajas con líneas más gruesas indican los protocolos destacados en esta guía.
\end{itemize}

\includegraphics[width=\linewidth]{protocol_dependencies.pdf}

\begin{exercise}
Queremos enviar por correo electrónico la frase ``Perdona, llegaré 5 minutos tarde''.
Nuestro cliente de correo enviará esos datos usando el protocolo SMTP, que a su vez
usa TCP (parte de la capa 4).
% 
Suponiendo que SMTP usa cabeceras en su \concept{PDU}, ¿cuántas cabeceras contendrá 
la \concept{PDU} correspondiente de la capa 2?
\end{exercise}

\begin{exercise}
¿Por qué crees que la capa~3 tiene un solo protocolo (dos si cuentas IPv4 e IPv6 por separado)? 
\end{exercise}

\section{¿Quién manda en Internet?}

El \conceptRef{IETF}{Internet Engineering Task Force (IETF)} se creó en 1986
para desarrollar y publicar estándares destinados a usarse en Internet. 
% 
Su producto principal son los documentos \conceptRef{RFC}{``Request for Comment'' (RFC)},
que contienen descripciones técnicas completas de esos estándares.


Estos RFC los crean grupos de trabajo de expertos, 
para que cualquiera los use libremente y de forma gratuita. 
% 
Cualquiera puede participar en estos grupos de trabajo, incluidos investigadores,
gobiernos y partes interesadas de la industria. 
% 
Las decisiones se toman cuando el 80-90\% de los participantes está de acuerdo, así que ninguna 
entidad controla por sí sola el proceso de estandarización.
% 
\begin{remark}
``Rechazamos a los reyes, a los presidentes y las votaciones. 
Creemos en el consenso aproximado y en el código que funciona'' 
--- Dave Clark (IETF), 1992.

Consulta el DOI \href{https://dx.doi.org/10.1109/MAHC.2006.42}{\underline{10.1109/MAHC.2006.42}} para 
una breve historia de lucha por el poder.
\end{remark}

Los estándares RFC son descriptivos, no prescriptivos. Esto significa que nadie vigila el uso 
correcto de estos estándares. En cambio, los fabricantes tienen el incentivo de cumplir los RFC 
para que sus productos sean compatibles con los demás. 
% 
El valor de los RFC (y de cualquier otro estándar) depende de que
los usuarios los adopten. Por este motivo, los RFC pueden actualizarse y \conceptRef{retire (standard)}{retirarse}.

Los protocolos TCP/IP prevalecen porque son flexibles y funcionan razonablemente bien en la mayoría de las situaciones, 
aunque no sean óptimos en prácticamente ninguna de ellas. 
Además, los protocolos TCP/IP son gratuitos, a diferencia de los protocolos \conceptRef{privative}{privativos}
que obligaban a comprar la solución de un fabricante concreto.
% 
\begin{remark}
Hoy cuesta imaginar un mundo \textit{no} 
dominado por la arquitectura TCP/IP, 
pero hasta finales de los 90 se comercializaron y mantuvieron pilas completamente distintas.
Algunos ejemplos son AppleTalk e IPX/SPX, de Apple y Novell, respectivamente.
\end{remark}

% https://www.icann.org/en/blogs/details/cheers-to-the-multistakeholder-community-30-9-2016-en

Algunos aspectos de TCP/IP, como los identificadores únicos 
(incluido \concept{DNS}, el Domain Name System)
y los números reservados, tienen que ser acordados por todos para que Internet funcione correctamente.
% 
La \concept{IANA} (Internet Assigned Numbers Authority) trabaja en coordinación con el \concept{IETF}
para definir los \conceptRef{RFC}{RFCs} adecuados.

\begin{remark}
Desde 2016, la IANA está gestionada por una organización multisectorial sin ánimo de lucro 
(\conceptRef{ICANN}{ICANN - Internet Corporation for Assigned Names and Numbers}).
% 
Desde 2004, la responsabilidad de los dominios regionales 
(\eg, \url{.cat}, \url{.fr}, \url{.es}) se transfiere a 
organizaciones de ámbito continental llamadas registros regionales de Internet (\conceptRef{RIR}{RIRs}).
\end{remark}



% # Based on https://www.youtube.com/watch?v=oIezCGjxV3A
% year,text
% 1962,Packet Switching
% 1969,ARPANET
% 1974,TCP/IP
% 1986,IETF
% 1991,WWW
% 1993,CIDR addressing
% 1998,Google
% ``The Internet doesn't do what the endpoints can do (\eg, reliability or compression).''


% - In February 1980, the Institute of Electrical and Electronics Engineers (IEEE) started project 802 to standardize local area networks (LAN).[16][27 (https://en.wikipedia.org/wiki/Ethernet)

% "protocols": "the first sheet of a volume" (on which contents and errata were written

% \begin{remark}
% Extra:
% \begin{itemize}
%     \item Video: https://www.youtube.com/watch?v=oIezCGjxV3A (part 1), https://www.youtube.com/watch?v=RN4gSBTANUY (part 2)
%     \item Monographic: https://cacm.acm.org/research/on-the-hourglass-model/
% \end{itemize}
% \end{remark}


% \begin{itemize}
% \item [\textbf{Layer 1}:] \textbf{Physical}
%     \begin{itemize}
%     \item[Features:] Ship \concept{digital} \otherBase{0}s and \otherBase{1}s through the real, \concept{analog} world.
%     \item[Offers:] Send and receive blocks of bytes.
%     \end{itemize}
% 
% \item [\textbf{Layer 2}:] \textbf{Network}\\
% 
% 
% \item [\textbf{Layer 3}:] \textbf{Internet}\\
% \item [\textbf{Layer 4}:] \textbf{Transport}\\
% \item [\textbf{Layer 7}:] \textbf{Application}\\
% \end{itemize}
% 
% \begin{remark}
% The \concept{OSI} model is a conceptual stack of layers that divides problems a bit differently from \concept{TCP/IP}. 
% TCP/IP was created on the go while the infrastructure of the Internet started its development, became its de-facto standard,
% and is used everywhere. The OSI model is useful to analyze and design communication protocols, but has limited practical usage.
% \end{remark}
% 
% 
% The network stack is normally part of the \concept{operating system} (OS).
% Only the \inlineCode{send} and \inlineCode{receive} functions of Layer 4 are exposed to the users of the OS.
% %
% This way, data transmitted by the \concept{network interface} are guaranteed to be properly formatted and follow
% the protocols of each stack layer.
% 
% \begin{remark}
% In order to control the lower layers --\eg, to transmit custom PDUs--, one needs to be the \texttt{root} user
% or give the \texttt{CAP\_NET\_RAW} kernel capability to a user program
% (or the equivalent in privative systems).
% \end{remark}
% 
% 
% \section{}\label{sec:stack_layers:organization}
% 
% - internet organization 
% bodies and stuff

